\section{The series in one place}\label{sec:series}

This section states the main results of the sixteen earlier papers and their eight memos, grouped by what they are about
and not by when they appeared. Nothing is reproved. Each statement carries its citation and its main hypotheses; the
cited theorems have further parameters (the number of tapes, the alphabet, the time constant), which we suppress.
\Cref{tab:series} at the end lists the papers in order. A reader who wants only the conclusion can skip to
\cref{sec:wall}.

\subsection{The question, passes and the coding formulation}\label{sec:first}

The first paper~\cite{FirstPaper2026} fixes the formal statement of~(T), shows that for linear time it agrees with the
simplest case of the formulation of Harvey and van der Hoeven~\cite[Proposition~11.1, a written proof]{FirstPaper2026},
and proves the following.
\begin{itemize}[leftmargin=1.5em]
\item \emph{Passes.} A transposer with at least one work tape makes $\Omega(\log d)$ passes, whatever it does with its
data: for oblivious machines with no premise on the running time, and for all machines that run in time polynomial
in~$L$~\cite[Theorems~5.9 and~5.12]{FirstPaper2026}. A memo removes the runtime premise for halting machines that are
correct on all inputs~\cite{PassesMemo2026}. This is the only logarithmic lower bound on a resource of the machine itself
that the series proves for all machines. It counts passes, not steps, and it is tight: the folklore algorithm uses
$O(\log d)$ passes.
\item \emph{Folding.} Along its popular traces, every limit of the rescaled work-head motion of a linear-time
transposer satisfies an isoperimetric inequality that forces the heads to fold back
(Theorem~M*,~\cite[Theorems~7.8 and~7.9]{FirstPaper2026}). The proof is a cut argument: the information that crosses a
cut of the time--space diagram is bounded by the states and the work-tape cells that enter across the cut.
\item \emph{Coding.} Every linear-time transposer yields a causal network code for complete exchange on the block
network of a time--space diagram of linear length, correct on a $2^{-O(d\log d)}$ fraction of the inputs, and on all
inputs if the machine is oblivious~\cite[Theorems~8.10 and~9.9]{FirstPaper2026}. The \emph{Tape Coding Conjecture} says
that for large $d$ no causal code solves complete exchange on such a fraction of the inputs on the block network of any
diagram of linear length; it implies that (T) fails~\cite[Conjecture~9.1, Theorem~9.2]{FirstPaper2026}. Its plain form,
for codes correct on all inputs, excludes oblivious linear-time transposers~\cite[Theorem~9.9]{FirstPaper2026}. On
these networks every routing has rate $O(1/\log d)$, so a solving code beats routing by a factor $\Omega(\log d)$, and
the conjecture is equivalent to saying that coding beats routing by only $o(\log d)$ on tape
networks~\cite[Theorems~8.8 and~9.4]{FirstPaper2026}.
\end{itemize}
The pass bound at polynomial running time, Theorem~M* and the reduction to codes survive when the machine is correct
on only a $2^{-O(d\log d)}$ fraction of the inputs. The same paper studies an abstract model of the one technique claimed
to beat $O(L\log d)$, recursion on XOR-mixed address frames: there every gadget costs at least its number of wires,
whereas linear time would need less~\cite[Theorem~12.4 and Proposition~12.9]{FirstPaper2026}; the correspondence between
the model and the preprint's recursion is informal. It also records why a proof must use the shape of tapes: a bound
$o(\log n)$ on coding gains for bounded-degree graphs in general would give superlinear circuit lower bounds for cyclic
shifts and multiplication~\cite[Proposition~11.2]{FirstPaper2026}.

The third paper~\cite{ThirdPaper2026} shows that the bit case is at least as hard: linear-time transposition of bit
squares of dyadic side implies~(T), by one fixed wrapper. The converse is open, and the paper proves barriers for
restricted compilers in that direction; one of them is a multiscale obstruction for compilers with boundedly many
stages.

\subsection{What head motion must do}\label{sec:motion}

A second line asks what the heads of a linear-time transposer must do.
\begin{itemize}[leftmargin=1.5em]
\item \emph{Not affine.} Along the popular traces, the limits of the rescaled work-head paths of a linear-time
transposer fail to be locally affine on sets of times of total measure at least $\tfrac12$, summed over the heads and
measured in blocks. Bounded-length words over three explicit palettes of scan operations are excluded, and so are
contracting recursions all of whose stages have regular limiting
motion~\cite[Theorems~6.3, 7.1 and~7.5]{SecondPaper2026}. No machine whose work heads are home-based with legs made of
boundedly many constant-rate runs transposes in linear time, and in every rescaled limit a linear-time transposer has
a linear amount of curved travel~\cite[Theorems~3.10 and~3.12]{FourthPaper2026}.
\item \emph{Macroscopic crossings.} An oblivious linear-time transposer routes a linear amount of information through
pairs of interleaved revisits whose excursions are a fixed fraction of~$L$: for every set of at least $\alpha L$ output
bits, every set of steps separating them from the earlier ends of such crossings has at least $c$ times as many
elements~\cite[Theorem~5.5]{SixteenthPaper2026}. A fixed fraction of this expansion sits on one coarsely monotone leg of
one tape, in a skeleton from a finite family independent of~$d$, and two work heads must each make linearly many long
revisits~\cite[Theorems~6.7 and~4.1]{SixteenthPaper2026}. The crossing and skeleton statements hold on heavy traces of
adaptive machines and of a bit simulator~\cite[Theorem~8.1]{SixteenthPaper2026}.
\item \emph{Rows to columns.} In the diagram of an oblivious transposer, every set of steps that separates the inputs
of a set of rows from the outputs they feed into a set of output rows carries the corresponding information, whatever
the coding~\cite[Lemma~7.1]{SixteenthPaper2026}. When the cuts for different address bits are disjoint this gives a
time bound $\Omega(L\log d)$ that no coding can evade~\cite[Corollary~7.3]{SixteenthPaper2026}.
\item \emph{Source covers.} For one fixed geometry of returns, a certificate that retains every input on which an
output can depend, and charges the information used to simplify the dependence, gives a correctness deficit of at
least $\eps L$ when the return families consist of a two-class core, at most one arbitrary exceptional family and
families of at most $\eps L$ hops, for fixed bounds on the running time, the capacity and the number of span
classes~\cite[Corollary~4.3]{SourceCovers2026}. A memo excludes funnels and single-sweep terminal landings and shows
where the certificate stops~\cite{GapAreaMemo2026}. These hypotheses on the geometry are not derived for arbitrary
machines.
\end{itemize}

Three results mark the limits of this line.
\begin{itemize}[leftmargin=1.5em]
\item \emph{Padding.} For every machine there is a wrapper, with more tapes and states and additive overhead
$O(d^2b)$, whose complete runs satisfy the cut inequality behind Theorem~M* with zero slack, and whose normalized head
limits satisfy all the conclusions of Theorem~M* under a linear time bound. The wrapper's motion is independent of the
computation that follows~\cite[Theorem~8.1]{SecondPaper2026}. It needs no idle tapes, and it also meets the finer
entry-cell form of the inequality at every slot length from one cell up to a constant number of
rows~\cite{PaddingMemo2026}. So these conditions on motion can be met by a machine that transposes nothing; they do not
see which information reaches the outputs.
\item \emph{One cut for all pairs.} In the diagram of any run, one set of steps separates the inputs from the outputs
for every pair of row sets at once~\cite[Proposition~10.2]{SixteenthPaper2026}. So counting the cuts of the
rows-to-columns inequality cannot by itself go beyond linear time: the disjointness that
\cite[Corollary~7.3]{SixteenthPaper2026} needs is not automatic.
\item \emph{The candidates.} Five explicit linear-time cascade schedules meet the crossing and skeleton conditions
at the sizes tested. All five are excluded, for fixed work alphabets, by the regular-recursion obstruction of the second
paper~\cite{CascadeObstructionMemo2026}, and two of them also by explicit cuts of the rows-to-columns
inequality~\cite[Theorem~10.8]{SixteenthPaper2026}. Whether the combined conditions on motion decide the oblivious case
is open.
\end{itemize}

\subsection{Depth reduction}\label{sec:depth}

The fourth paper~\cite{FourthPaper2026} splits the oblivious case of~(T) into two statements.
\begin{itemize}[leftmargin=1.5em]
\item \emph{Depth reduction} (the trajectory cut problem $\TCP(k)$): for every $\eps>0$, deleting $\eps T$ edges from
the dependency graph of $k$ head walks of length~$T$ leaves no path with more than a bounded number of long hops, and
then, after periodic steps are also revealed, paths of bounded length.
\item \emph{The crux}: an index-coding lower bound on the resulting bounded cones, strong enough to exceed the cost of
the revealed steps.
\end{itemize}
Together they would exclude oblivious linear-time
transposers~\cite[Theorem~4.10 and Corollary~4.11]{FourthPaper2026}. Depth reduction is now a theorem. It was proved
for home-based walks~\cite[Theorem~4.20]{FourthPaper2026}; then for walks of bounded nesting load, and for one arbitrary
head together with heads of bounded nesting load at boundedly many
scales~\cite[Theorems~3.6 and~4.2]{EighthPaper2026}; and finally for every number of walks, with a depth $D(k,\eps)$
such that $\log\log D=O(k^2\log(k/\eps))$~\cite[Theorem~7.5]{BoundedStack2026}. A memo shows that $\TCP(k)$ for all~$k$
is equivalent to constant-depth revelation of arbitrarily small entropy per input bit, for all linear-time oblivious
computations~\cite{InfoDepthMemo2026}. The fourteenth paper also refutes Conjectures~A$'$ and~A$''$ of the fourth and
eighth papers, which would have made the deep-nesting pairs
depth-robust~\cite[Corollary~9.1]{BoundedStack2026}: walks nested at $m$ scales are cut to a depth independent of~$m$.
It draws the conclusion that matters here~\cite[Corollary~9.3]{BoundedStack2026}:
\begin{quote}
\emph{for oblivious machines, the crux alone would exclude linear-time transposers.}
\end{quote}
In the language of \cref{sec:storage}: for an oblivious linear-time transposer and every $\eps>0$, revealing $\eps L$
bits at suitable steps leaves every output bit a function of the revealed word and of boundedly many input bits, and
for some value of the revealed word the inputs compatible with it form a binary storage code of rate at least $1-\eps$
on the machine's own cone relation. The same genies give depth reduction at the level of
codes~\cite[Corollary~9.6]{BoundedStack2026}, as conjectured in the eleventh paper. And, given the \emph{native half
bound} (every binary storage code on the cone relations of the machine's own traces has rate at most
$\tfrac12+o(1)$, which is open), they exclude every linear-time transposer, adaptive or not, whose trace has entropy at
most $\eta L$ for a fixed $\eta<\tfrac12$~\cite[Corollary~9.7]{BoundedStack2026}.

For machines whose data path is linear over a field, linear depth reduction together with a \emph{linear crux}, a
min-rank bound for the ball systems of schedules, excludes oblivious linear-time
transposers~\cite[Theorem~7.5]{EighthPaper2026}. The linear crux is exactly the relative rigidity of the transposition
matrix described in \cref{sec:onehop}~\cite[Problem~10.6]{EighthPaper2026}, and since the fourteenth paper it alone
suffices~\cite[Corollary~9.3(c)]{BoundedStack2026}.

\subsection{Storage codes: what is proved}\label{sec:proved}

The crux needs an upper bound on the rate of storage codes on cone relations. The series proved such bounds, at the
pairing rate, for many classes of relations.
\begin{itemize}[leftmargin=1.5em]
\item \emph{Homomorphisms and small degree.} Rate bounds pull back along homomorphisms, so every storage code on a
graph that maps to $C_5$, to $K_{14/5}$ or to the Petersen graph has rate at most
$\tfrac12$~\cite[Theorem~6.5 and Corollary~6.9]{FourthPaper2026}. On triangle-free graphs of maximum degree three no
scalar linear code over any field beats~$\tfrac12$~\cite[Theorem~6.13]{FourthPaper2026}. The Clebsch graph, of degree
five, has guessing number~$10$ on $16$ vertices~\cite[Corollary~6.11]{FourthPaper2026}, so bounded degree alone is not
enough.
\item \emph{Bounded complex coefficients.} On locally bipartite graphs of bounded degree, linear codes with bounded
complex coefficients have rate at most $\tfrac12+o(1)$~\cite[Theorem~6.17]{FourthPaper2026}. The same holds for complex
linear recovery with bounded coefficients and fixed block width on directed relations of bounded incidence degree that
are separated, hence on cone relations after a paid trimming of
degrees~\cite[Theorem~4.5 and Corollary~4.7]{DirectedRecovery2026}. On such relations a gain by linear codes with
boundedly many symbols per vertex therefore needs positive characteristic, or complex coefficients that no diagonal
rescaling makes bounded.
\item \emph{Hyperfinite relations.} Every code over every finite alphabet on relations whose incidence graphs are
uniformly hyperfinite and separated has rate at most $\tfrac12+o(1)$; no degree bound is
needed~\cite[Corollary~4.4]{DirectedRecovery2026}.
\item \emph{Nested elimination.} The exact bound $\card C^2\le\card\alpha^{\card V}$ holds for all codes, linear or not,
on large classes of relations without transitive triangles: triangle-free graphs of maximum degree three; loopless
relations whose underlying graph has maximum degree three; relations whose vertices of in-degree at least three induce
an acyclic sub-relation; relations partitioned into blocks with private outside parents, each block chordless
subcubic, a chorded digon path or internally certified; and Cayley relations $\mathrm{Cay}(G,\{s,s^{-1},t\})$ of every
finite group~\cite[Theorems~4.3, 5.2, 5.3 and~5.10, Corollary~5.12]{SeventhPaper2026}. It also holds on the class of
the fourth paper's Theorem~S, for all codes~\cite[Theorem~3.6]{DirectedReciprocity2026}. The degree bound cannot be
raised: a $4$-regular triangle-free graph on $13$ vertices carries a linear code over $\F_4$ with
$\card C^2=4^{14}$~\cite[Example~7.8]{SeventhPaper2026}.
\item \emph{Covers and towers.} For the binary adjacency code of the Clebsch graph, every connected tower of binary
lifts has kernel dimension at most $N/2+2$~\cite[Theorem~9.2]{Covers2026}; lifting a binary linear code with a
unit-diagonal check along a tower of double covers never raises its rate~\cite[Theorem~8.1]{EighthPaper2026}.
\item \emph{Memos.} The half bound for circular cliques $K_{p/q}$ extends to every $p/q<3$~\cite{CircularMemo2026};
nested elimination extends to soft recovery with record capacities~\cite{SoftRecovery2026}; and the bounds for supports
of maximum degree three extend beyond scalar linear codes, to reversible rules over any alphabet and to linear codes
whose rows have at most one port of rank above one~\cite{SubcubicMemo2026}.
\end{itemize}

One result of this group is about the relations of machines themselves. For cones of fixed depth in linear-time
schedules, after a paid trimming of degrees, only $O(L/\log d)$ matched pairs of an input port and the output that
demands it lie within any fixed incidence
distance~\cite[Theorem~3.3, Corollaries~3.7 and~3.8]{DirectedRecovery2026}. So cone relations are separated in the
sense of \cref{sec:storage}, and almost no output is within bounded distance of its own input bit. No gadget of bounded
size can then lift a code above the pairing rate: a gain above $\tfrac12$ on boundedly many coordinates needs a short
path between the two ports of one of them~\cite[Lemma~4.1 and Theorem~4.2]{DirectedRecovery2026}. The $L/\log d$ comes
from a bound in steps for transport at bounded occupancy, which charges every transported label the $\log d$ bits of
its row address~\cite[Lemma~3.4]{DirectedRecovery2026}: for uncoded movement the logarithm is proved, and what the
bound does not cover is a cell that carries a combination of many labels.

This led to the \emph{all-code question}~\cite[Question~9.1]{DirectedRecovery2026}: does every storage code on
separated relations of bounded incidence degree have rate at most $\tfrac12+o(1)$?

\subsection{Storage codes: what is false}\label{sec:false}

The answer is no.
\begin{itemize}[leftmargin=1.5em]
\item \emph{Above the pairing rate.} There are relations of bounded (astronomically large) incidence degree in which
every paired distance tends to infinity, with $\F_{127^2}$-linear storage codes of rate above $0.5021$, and binary
codes of rate above $0.5013$; for large primes inert in $\Q(\sqrt5)$ the rates approach
$4-2\sqrt3\approx0.536$~\cite[Theorems~7.1 and~7.3, Corollary~8.2]{StorageCodes2026}. The relations are twisted
quotients of a Bruhat--Tits tree, the codes are Tanner codes with Reed--Solomon local codes, and their dimension comes
from Riemann--Roch on Hilbert modular surfaces in characteristic~$p$. The undirected form of the question, on graphs of
bounded degree with long odd cycles, has a negative answer as well~\cite[Corollary~9.3]{StorageCodes2026}. Tanner operators of the same
kind, at $p=3$, refute the independence of the chain in Jaikin-Zapirain's arbitrary-field form of the L\"uck
approximation conjecture~\cite[Theorem~5.6]{Luck2026}.
\item \emph{Rates approaching one.} A linear-time transposer, even an adaptive one, would yield binary storage codes
of this kind at every large degree cap~$\Delta$, with rates tending to one as $\Delta$
grows~\cite[Theorem~3.1]{HighRate2026}; and such codes exist abstractly, with rate at least
$1-C\Delta^{-\alpha}$~\cite[Theorem~5.5]{HighRate2026}. So, at the scale of that transfer theorem, no rate bound for
separated relations of bounded degree in general can exclude a transposer~\cite[Corollary~6.2]{HighRate2026}: a lower
bound must use a further restriction imposed by the physical computation.
\end{itemize}
No tape process is known that realizes the recoveries of either family.

\subsection{Adaptive machines}\label{sec:adaptive}

Much of the above already covers adaptive machines: the pass bound (at polynomial running time, or by the memo),
Theorem~M* along popular traces, and the reduction of~(T) to the Tape Coding Conjecture. Beyond that the series has the
following.
\begin{itemize}[leftmargin=1.5em]
\item Every linear-time transposer yields an external-memory transposition with $O(d)$ I/Os on blocks of one row that
is oblivious on $2^{L-O(d)}$ inputs, so an external-memory lower bound without the indivisibility assumption would
refute~(T)~\cite[Theorem~6.4 and Corollary~6.5]{MergeDepth2026}.
\item The source-cover exclusions extend to machines of sufficiently low
adaptivity~\cite[Theorem~6.1]{MergeDepth2026}.
\item Given the native half bound for the relations of skeleton classes, which is open, a linear-time transposer whose
work heads enter every cell boundedly often must carry about $L/2$ bits in its coarse skeleton, hence in the coarse
motion of its input head~\cite[Theorem~4.13 and Corollary~4.16]{MergeDepth2026}; the fourteenth paper removes the
hypothesis on the work heads from the skeleton statement~\cite[Corollary~9.7]{BoundedStack2026}.
\item The fifteenth paper extracts its codes from any linear-time transposer through a packet
compilation~\cite[Theorem~3.1]{HighRate2026} and, at the original bit positions, from the traces of a bit simulator
with few traces~\cite[Theorems~8.1 and~9.2]{HighRate2026}. The crossing theorems hold on heavy traces of the machine and
of that simulator~\cite[Theorem~8.1]{SixteenthPaper2026}; whether they hold in the machine's own diagram on traces that
carry almost all inputs was left open there~\cite[Problem~11.2]{SixteenthPaper2026}. A paper that followed this one
proves a statement of that form at the canonical entry width: on the original run of all but an exponentially small
share of the inputs that a linear-time transposer handles, every cut between the outputs of a large set and the
macroscopic crossings has size proportional to that set~\cite[Theorem~H.7]{AdaptiveTransfer2026}. For every entry
width, and again outside an exponentially small share of the inputs, it proves that two work heads make linearly many
macroscopic returns, and that cuts between a large set of outputs and the reads of their inputs are of linear
size~\cite[Theorem~4.3 and Corollary~5.4]{AdaptiveTransfer2026}.
\end{itemize}

\begin{table}[t]
\centering\small
\begin{tabular}{@{}r>{\raggedright\arraybackslash}p{0.27\textwidth}>{\raggedright\arraybackslash}p{0.62\textwidth}@{}}
\toprule
 & Paper & Main contribution to the question \\
\midrule
1 & Tape Coding Conjecture~\cite{FirstPaper2026} & Formal statement; $\Omega(\log d)$ passes; Theorem~M*; transposers
are network codes \\
2 & Local Coding Barriers~\cite{SecondPaper2026} & Local coding bounds; affine and regular motion excluded; padding \\
3 & Bit and Word~\cite{ThirdPaper2026} & Transposition of dyadic bit squares implies~(T) \\
4 & Curved Motion~\cite{FourthPaper2026} & Curved motion; depth reduction plus crux; storage codes at rate $\tfrac12$ \\
5 & Directed Recovery~\cite{DirectedRecovery2026} & Paired-distance separation of cones; half bounds; the all-code
question \\
6 & Covers~\cite{Covers2026} & Half bounds along covers and for locally generated codes \\
7 & Nested Elimination~\cite{SeventhPaper2026} & Exact all-code half bounds by elimination \\
8 & Bounded Nesting~\cite{EighthPaper2026} & Depth reduction at bounded nesting; the linear crux as rigidity \\
9 & Source Covers~\cite{SourceCovers2026} & Source-cover certificates and exclusions \\
10 & Directed Reciprocity~\cite{DirectedReciprocity2026} & All-code bound on the class of Theorem~S; deep nesting \\
11 & Merge Depth~\cite{MergeDepth2026} & Merge depth; adaptive machines; external memory \\
12 & L\"uck Approximation~\cite{Luck2026} & A Tanner operator whose kernel survives along a congruence tower and dies
along pro-$3$ towers \\
13 & Above Half Rate~\cite{StorageCodes2026} & The all-code question has a negative answer \\
14 & Bounded-Stack DAGs~\cite{BoundedStack2026} & Depth reduction for every number of walks \\
15 & High-Rate Codes~\cite{HighRate2026} & A transposer needs rates tending to one; such codes exist abstractly \\
16 & Macroscopic Crossings~\cite{SixteenthPaper2026} & Crossings, skeletons, rows to columns, one cut for all pairs \\
\bottomrule
\end{tabular}
\caption{The earlier papers of the series.}\label{tab:series}
\end{table}
